% ECONOMICS_PROBLEM: {"id": "L41-501", "title": "Merger efficiencies", "jel_code": "L41", "source_id": "OP-0394"}
% !TeX program = xelatex
\documentclass[11pt,a4paper]{article}
\usepackage[margin=23mm]{geometry}
\usepackage{fontspec}
\setmainfont{Latin Modern Roman}
\usepackage{amsmath,amssymb,mathtools}
\usepackage{xcolor,enumitem,booktabs,longtable,array,xurl,hyperref}
\definecolor{muted}{HTML}{53636C}
\hypersetup{colorlinks=true,urlcolor=blue,linkcolor=blue}
\setlength{\parindent}{0pt}
\setlength{\parskip}{5pt}
\newcommand{\R}{\mathbb R}
\newcommand{\E}{\mathbb E}
\newcommand{\N}{\mathbb N}
\newcommand{\ind}{\mathbf 1}
\newcommand{\Var}{\operatorname{Var}}
\newcommand{\Cov}{\operatorname{Cov}}
\newcommand{\supp}{\operatorname{supp}}
\DeclareMathOperator*{\argmin}{arg\,min}
\DeclareMathOperator*{\argmax}{arg\,max}
\newcommand{\entrymeta}[3]{{\sffamily\small\color{muted}\textbf{\texttt{#1}}\quad|\quad #2\par #3\par}}
\newcommand{\Problemref}[1]{\texttt{#1}}
\newcommand{\JELref}[2]{\texttt{#2} (JEL \texttt{#1})}
\begin{document}
\section*{Merger efficiencies}
\textbf{Problem ID:} \texttt{L41-501}\quad \textbf{JEL:} \texttt{L41}\par
% BEGIN ECONOMICS STATEMENT
\label{problem:OP-0394}
{\sffamily\small\color{muted}Original subject group: Industrial organization and competition\par}
\label{definitions:problem:30007676}
\textbf{Audit classification:} Published research agenda. The original claimed marginal-cost savings and realized total-cost difference could have incompatible units. The joint distribution of individual causal effects is generally not identified from separate randomized regime marginals.
\subsection*{Definitions and precise research target}
Consider proposed horizontal manufacturing mergers reviewed in the United States over a prespecified five-year cohort. For deal \(m\), let \(S_m\) be claimed annual real-cost savings at the fixed reference output documented before approval, and \(X_m\) other premerger information. Let \(C_m(a,h)\) denote the cost of producing a fixed reference output vector at horizon \(h=3\) years under merger regime \(a\in\{0,1\}\), and let \(P_m(a,h)\) be the corresponding expenditure-weighted consumer price index after equilibrium repricing. Define realized cost savings \(G_m=C_m(0,h)-C_m(1,h)\), the claim-calibration function \(g(s,x)=E[G_m\mid S_m=s,X_m=x]\), and consumer price benefit \(B_m=P_m(0,h)-P_m(1,h)\). Costs and prices are measured in constant dollars; expectation refers to the reviewed-deal population. Obtain production and logistics audits that distinguish technological savings from purchasing-power transfers and output changes. Use approved and appropriately comparable abandoned or blocked deals to construct counterfactuals, recording the selection assumptions needed for each comparison. Identify or bound, using an explicitly stated coupling restriction when needed, the joint distribution of \(G_m\) and \(B_m\), rather than attributing all profit growth to efficiency. Determine which preapproval evidence predicts positive, merger-specific savings and consumer benefits. Where approval selection precludes identification, provide bounds and specify the additional information that would reduce them.

\textbf{Definitions, assumptions and mathematical qualifications.} Choose a reference annual output vector \(\bar q_m\) and a constant-dollar base year before observing merger outcomes. \(C_m(a,h)\) is annual real resource cost of producing that vector at horizon \(h\), including an explicit annual allocation of fixed production resources. \(S_m\) must be the claimed annual cost reduction for that same vector; a change in cost per unit is a different object. \(P_m(a,h)\) is a fixed-basket price level with stated base-year expenditure weights. \(G_m\) and \(B_m\) are contrasts of potential cost and price outcomes under a common equilibrium selection rule. Purchasing-power transfers are distinct from technological resource savings. Conditional expectations \(g(s,x)\) are defined only almost everywhere on the support of \((S_m,X_m)\).
\subsection*{Variants and assumption changes}
\begin{enumerate}
\item \textbf{Editorial, calibration.} Study \(E[G_m-S_m\mid X_m]\) and \(E[B_m\mid G_m,X_m]\) under a design identifying each contrast; observing merged-firm profits alone identifies neither.
\item \textbf{Editorial, endogenous output.} Replace fixed-output savings by total surplus under each regime, allowing demand, output, investment and entry responses; separate this outcome from engineering efficiency at fixed output.
\end{enumerate}
\subsection*{References and source audit}
\textbf{Checked source.} 24 August 2024 author draft, Section 6, printed p. 37, sixth research direction. Full primary text retrieved. \url{https://web.stanford.edu/~ayurukog/trendsincompetition.pdf}.
\begin{enumerate}\raggedright
\item\hypertarget{definitions:source:30007676:1}{} Carl Shapiro; Ali Yurukoglu. 2024. \emph{Trends in Competition in the United States: What Does the Evidence Show?}. 24 August 2024 author draft, Section 6, printed p. 37, sixth research direction.
\url{https://web.stanford.edu/~ayurukog/trendsincompetition.pdf}.
DOI: \url{https://doi.org/10.3386/w32762}.
\end{enumerate}

\subsection*{Common conventions: Source mathematical conventions}

These conventions apply to each entry unless explicitly replaced.

\textbf{Models and identification.} A primitive model \(m\) specifies all probability laws, feasible choices, information, equilibrium and measurement rules used by its target. The admissible class \(\mathcal M\) is fixed independently of outcomes. If \(P_O(m)\) is its observable probability law and \(\tau(m)\) the target, its identified set at observable law \(P\) is
\[
 \mathcal I_\tau(P)=\{\tau(m):m\in\mathcal M,\ P_O(m)=P\}.
\]
Point identification means that this set is a singleton. An empty set signals incompatibility of data and assumptions. Coordinate bounds need not describe the joint set. Bounds are sharp only if every retained target is attainable or arbitrarily approachable by compatible models. Finite bounds require explicit restrictions on unobserved quantities; unrestricted confounding, hidden wealth, extrapolation or arbitrary equilibrium selection can yield uninformative sets.

\textbf{Causal interventions.} A potential outcome \(Y_i(p)\) is the outcome under a complete policy regime \(p\), including timing, financing, information and market spillovers. The target population and units are fixed before treatment. With interference, use the complete assignment vector or a justified exposure mapping; individual no-interference notation alone cannot identify an economy-wide intervention. A difference of mean potential outcomes does not require the cross-world joint distribution. Conditioning on post-treatment migration, survival, enrollment or employment changes the estimand and needs separate justification.

\textbf{Statistical statements.} An estimator, confidence set, forecast or test specifies a sampling law, model class and loss/coverage criterion. Uniform \((1-\alpha)\)-coverage means \(\inf_{m\in\mathcal M}\Pr_m\{\tau(m)\in C_n\}\geq1-\alpha\), or an explicitly asymptotic version. The whole parameter space trivially has coverage, so informativeness requires a stated width, risk or power objective. A forecasting improvement is relative to a named benchmark, information set and evaluation population.

\textbf{Equilibrium and optimization.} An equilibrium comprises feasible plans, optimality, beliefs where needed, budget constraints and all market/resource clearing. The primitives or a stated benchmark define each correspondence \(\mathcal E(p;m)\). Nonemptiness is assumed or proved, never inferred just from compact policy sets. A selected-equilibrium target uses a declared selection rule; a robust target quantifies over all permitted equilibria. Use suprema or infima unless attainment is established. An \(\varepsilon\)-optimal policy is feasible and within \(\varepsilon\) of the finite optimum value. Report an empty feasible set separately.

\textbf{Welfare and accounting.} Normative utility, interpersonal normalization, social weights, numeraire, discounting and the use of public receipts are primitives. Behavioral utility need not coincide with the welfare criterion. Household welfare includes ownership income and rents once; adding profits again to compensating variation generally double-counts them. A monetary equivalent is defined by an explicit transfer and price convention, with monotonicity and range conditions. Average effects, mechanism contributions, transfers and welfare losses are different targets. Infinite sums require convergence and dynamic feasible plans require terminal or no-Ponzi conditions.

\textbf{Source versions.} A working-paper date, revision date and journal publication year are distinguished. A DOI for a journal version is not automatically the DOI of an earlier manuscript. Locators refer to the stated version and printed pages where specified. A metadata-only check does not verify a claim about a full-text conclusion. Every entry's source subsection narrows the provenance claim accordingly.
% END ECONOMICS STATEMENT
\end{document}
